\documentclass[aspectratio=169]{beamer}
\usetheme{metropolis}
\usepackage{appendixnumberbeamer}
\usepackage{booktabs, hyperref}

\input{mgmt675-style}

\usepackage{tikz}
\usetikzlibrary{shapes.geometric, arrows.meta, positioning, calc}

% Title info
\subtitle{MGMT 675: Generative AI for Finance}
\title{Module 2: Portfolio Optimization}
\author{Kerry Back}
\date{}

\begin{document}

\maketitle

\begin{frame}{Learning Objectives}
\begin{enumerate}
  \item Construct optimal portfolios with real-world constraints using AI
  \item Visualize the opportunity set, efficient frontier, and capital allocation line
  \item Understand how constraints change the attainable risk-return tradeoff
\end{enumerate}

\vspace{0.5cm}
These are core tools of investment management.  AI handles the mathematics and code.  \alert{You focus on the assumptions and the interpretation.}
\end{frame}

% ============================================================
% PORTFOLIO OPTIMIZATION
% ============================================================
\section{Portfolio Optimization}

\begin{frame}{The Problem}
Given expected returns, risks (standard deviations), and correlations for multiple assets, plus a risk-free rate:

\begin{itemize}
  \item What is the \textbf{best portfolio}?
  \item What does the set of all possible portfolios look like?
  \item How do real-world constraints change the answer?
\end{itemize}

\vspace{0.3cm}
Mean-variance analysis provides the framework.  AI handles the computation.
\end{frame}

\begin{frame}{Key Concepts}
\begin{columns}[T]
\begin{column}{0.48\textwidth}
\begin{shadedbox}[title=\textbf{Portfolios}]
\begin{itemize}\small
  \item \textbf{Tangency portfolio:} the risky portfolio with the highest Sharpe ratio (excess return per unit of risk)
  \item \textbf{Global minimum variance (GMV):} the risky portfolio with the lowest possible risk
  \item \textbf{Efficient frontier:} the set of risky portfolios with the highest return for each level of risk
\end{itemize}
\end{shadedbox}
\end{column}
\begin{column}{0.48\textwidth}
\begin{shadedbox}[title=\textbf{Capital Allocation Line}]
\begin{itemize}\small
  \item The straight line from the risk-free rate through the tangency portfolio
  \item Represents all combinations of the risk-free asset and the tangency portfolio
  \item The best attainable risk-return tradeoff
\end{itemize}
\end{shadedbox}
\end{column}
\end{columns}
\end{frame}

\begin{frame}{Visualizing the Opportunity Set}
\begin{center}
\begin{tikzpicture}[scale=0.7]
  % Axes
  \draw[->, thick] (0,0) -- (10,0) node[right] {Standard Deviation};
  \draw[->, thick] (0,0) -- (0,6.5) node[above] {Expected Return};

  % Portfolio cloud (scatter dots)
  \foreach \i in {1,...,60}{
    \pgfmathsetmacro{\x}{1.5 + 6*rnd}
    \pgfmathsetmacro{\y}{0.5 + 5*rnd}
    \pgfmathsetmacro{\valid}{\y < 0.4*\x + 2.5 ? 1 : 0}
    \ifnum\valid=1
      \fill[accentblue!20] (\x,\y) circle (1.5pt);
    \fi
  }

  % Efficient frontier (curve)
  \draw[very thick, alertorange, domain=1.8:8, samples=50]
    plot (\x, {0.7*sqrt(\x-1) + 1.5});

  % CAL (line from risk-free through tangency)
  \draw[very thick, accentblue, dashed] (0,1.2) -- (9, 5.7);

  % Risk-free rate
  \fill[titlegray] (0,1.2) circle (3pt) node[left, font=\small] {$r_f$};

  % Tangency portfolio
  \fill[alertorange] (4.5,4.0) circle (4pt) node[above right, font=\small\bfseries] {Tangency};

  % GMV
  \fill[titlegray] (1.8,2.3) circle (4pt) node[below right, font=\small\bfseries] {GMV};
\end{tikzpicture}
\end{center}

\small The \textcolor{accentblue!50}{cloud} shows random portfolios.  The \textcolor{alertorange}{orange curve} is the efficient frontier.  The \textcolor{accentblue}{dashed line} is the capital allocation line.
\end{frame}

\begin{frame}{Adding Real-World Constraints}
\begin{columns}[T]
\begin{column}{0.48\textwidth}
\textbf{Common constraints}
\begin{itemize}
  \item No short sales (all weights $\geq 0$)
  \item Maximum position (e.g., $\leq 40\%$ per asset)
  \item Minimum position (e.g., $\geq 2\%$ if held)
  \item Margin requirement (sum of absolute values $\leq 2$)
\end{itemize}
\end{column}
\begin{column}{0.48\textwidth}
\textbf{What happens}
\begin{itemize}
  \item The efficient frontier shifts \alert{inward} (lower returns and/or higher risk)
  \item The tangency portfolio changes
  \item The Sharpe ratio cannot improve with constraints---it can only stay the same or get worse
\end{itemize}
\end{column}
\end{columns}
\vspace{0.3cm}
\alert{The constrained frontier is still useful: it reflects what you can actually implement.}
\end{frame}

\begin{frame}{Solution Methods}
\begin{itemize}
  \item \textbf{Solver (numerical optimization)}
    \begin{itemize}
      \item Maximize Sharpe ratio (tangency portfolio)
      \item Minimize risk for a target return (efficient frontier)
      \item Minimize risk (global minimum variance portfolio)
    \end{itemize}
  \item \textbf{Analytic/algebraic solution}
    \begin{itemize}
      \item Solve systems of linear equations for tangency, GMV, and frontier
      \item Available only when there are no constraints
    \end{itemize}
\end{itemize}

\vspace{0.3cm}
Python solver options: \texttt{scipy.minimize}, \texttt{cvxopt}, \texttt{cvxpy}---AI chooses and writes the code.

\vfill
\alert{Ask Claude:} \small Discuss the advantages and disadvantages of these solver options for mean-variance analysis with inequality constraints.
\end{frame}

\begin{frame}[shrink=15]{Exercise: Portfolio Cloud}

Consider three assets with the following characteristics:

\vspace{0.2cm}
\begin{itemize}\small
  \item \textbf{Asset A}: expected return 8\%, standard deviation 14\%
  \item \textbf{Asset B}: expected return 12\%, standard deviation 20\%
  \item \textbf{Asset C}: expected return 15\%, standard deviation 26\%
  \item Correlations: $\rho_{AB} = 0.3$, \quad $\rho_{AC} = 0.1$, \quad $\rho_{BC} = 0.5$
  \item Risk-free rate: 4\%
\end{itemize}

\vspace{0.2cm}
Ask Claude to:
\begin{enumerate}
  \item Generate 10{,}000 random portfolios (random weights summing to 1, no short sales) and plot each portfolio's expected return vs.\ standard deviation.
  \item Overlay the efficient frontier and the capital allocation line on the same plot.
  \item Mark the tangency portfolio and the global minimum variance portfolio.
\end{enumerate}
\end{frame}

% ============================================================
% EXERCISES
% ============================================================
\section{Exercises}

\begin{frame}{Exercise 1: Tangency Portfolio}
\begin{itemize}
  \item Download data for 5 ETFs + T-bill rate; compute tangency portfolio weights, return, standard deviation, and Sharpe ratio.
  \item Plot the efficient frontier and capital allocation line, then add a no-short-sales constraint and compare.
  \item Submit the data, plots, and a short paragraph on how constraints changed the result.
\end{itemize}
\end{frame}

\begin{frame}{Exercise 2: Portfolio Cloud Artifact}
Using the three-asset data from earlier:
\begin{itemize}\small
  \item \textbf{Asset A}: expected return 8\%, standard deviation 14\%
  \item \textbf{Asset B}: expected return 12\%, standard deviation 20\%
  \item \textbf{Asset C}: expected return 15\%, standard deviation 26\%
  \item Correlations: $\rho_{AB} = 0.3$, \quad $\rho_{AC} = 0.1$, \quad $\rho_{BC} = 0.5$; \quad $r_f = 4\%$
\end{itemize}
\vspace{0.2cm}
Ask Claude to build an \textbf{interactive artifact} that generates 10{,}000 random portfolios, overlays the efficient frontier, and marks the tangency and GMV portfolios.  Publish and submit the link.
\end{frame}

\end{document}
